\appendices
\section{Proofs of auxiliary results}
\label{app:proofs}

\subsection{Proof of Corollary~\ref{cor:deterministic-band}}
\begin{proof}
Under \eqref{eq:boundedconditional}, $R_G=\log \LR_{G\mid C}$ lies in $[\log m,\log M]$ almost surely. If $\LR_C<\lambda/M$, even the largest genomic factor cannot cross the threshold; if $\LR_C>\lambda/m$, even the smallest genomic factor cannot move the observation below threshold. This gives \eqref{eq:boundaryband}--\eqref{eq:logboundaryband}.
\end{proof}

\subsection{Proof of Corollary~\ref{cor:zebra}}
\begin{proof}
Equation \eqref{eq:zlr} shows that posterior odds differ from $\LR_C$ only by the positive constant prior odds $\pi/(1-\pi)$. Substituting into \eqref{eq:boundaryband} yields the result.
\end{proof}

\subsection{Proof of Theorem~\ref{thm:tail-extension}}
\begin{proof}
Consider the upper disease-evidence event $\{\log\LR_{C_n}\ge K\}$ with $K>0$. On
\begin{equation}
\{|R_{G,n}|\le b_\delta\}\cap\{\log\LR_{C_n}\ge K>0\},
\end{equation}
we have
\begin{equation}
\left|
\frac{\log\LR_{C_n,G}-\log\LR_{C_n}}
     {\log\LR_{C_n}}
\right|
=\frac{|R_{G,n}|}{\log\LR_{C_n}}
\le \frac{b_\delta}{K}.
\end{equation}
Therefore violation of the relative bound on this upper tail implies $|R_{G,n}|>b_\delta$, whose $\Pzero$ probability is at most $\delta$ uniformly in $n$ by A1-U. A2 is used only to ensure that positive-probability upper-tail events exist for arbitrarily large $K$; it is not derived from the preceding algebra.
\end{proof}

\subsection{Proof of Proposition~\ref{prop:genomicpower}}
\begin{proof}
For part (i), using the Radon--Nikodym identity under $\Pzero$,
\begin{align}
\Eone[\phi(G)]
&=\Ezero[\LR_G(G)\phi(G)]\\
&\le M_G\Ezero[\phi(G)]\\
&=M_G\alpha.
\end{align}
For part (ii), define $A_K=\{c:\LR_C(c)\ge K\}$. Then
\begin{align}
\Pone(A_K)
&=\Ezero[\LR_C(C)\ind\{A_K\}]\\
&\ge K\Pzero(A_K)=K\alpha_K,
\end{align}
which gives \eqref{eq:clinicalLRplus}. If $K>M_G$, part (i) implies that any genomic-only test with false-positive rate $\alpha_K$ has power at most $M_G\alpha_K<K\alpha_K\le \Pone(A_K)$. Markov's inequality applied to $\Ezero[\LR_C]=1$ yields $\alpha_K\le 1/K$.
\end{proof}

\subsection{Proof of Corollary~\ref{cor:genomic-postprocessing}}
\begin{proof}
Under the class-independent channel $Q$, the score likelihood ratio satisfies
\begin{equation}
\LR_S(s)=\mathbb{E}_0[\LR_G(G)\mid S=s].
\end{equation}
A conditional expectation of a random variable bounded above by $M_G$ is itself bounded above by $M_G$. For a stochastic learning algorithm, this statement is interpreted conditional on the fitted model (or training sample); the deployment-time map from a new subject's $G$ to $S$ must satisfy the stated class-independence condition.
\end{proof}

\subsection{Proof of Proposition~\ref{prop:auc-local}}
\begin{proof}
AUC is the probability that a randomly drawn case receives a higher score than a randomly drawn control, with the usual half-credit convention for ties. If neither member of a case--control pair lies in $B$, then $S'=S$ for both members and the pairwise ranking contribution is unchanged. Therefore the AUC contribution can differ only when at least one member of the pair lies in $B$. The probability of that event is
\begin{multline}
1-(1-\Pone(B))(1-\Pzero(B))\\
=\Pone(B)+\Pzero(B)-\Pone(B)\Pzero(B),
\end{multline}
which bounds the maximum possible AUC change.
\end{proof}
